\chapter{Carnap: comparative confirmation, relevance, and reliability}
\label{sec:p15-carnap-confirmation-relevance-and-reliability}

This chapter continues the first-edition 1950 presentation, not the later
continuum of inductive methods. Its sources are the original chapters
on relevance, comparative inductive logic, and estimation, including
their restrictions and countercases
\cite[secs.~65--76, 79--88, 98--107]{Carnap1950Probability}.
The state descriptions, sentence ranges, regular measures, and
\(c(h,e)=m(e\land h)/m(e)\) are those of
Chapter~\ref{sec:p15-carnap-logical-probability-and-induction}.
All finite-system statements allow arbitrary sentences of the language,
including quantified sentences: quantifiers are evaluated over its entire
named finite domain. No particular \(c^*\), monadic vocabulary, or
empirical frequency rule is imposed unless stated.

\section{Comparison without selecting one measure}
\label{sec:p15-carnap-comparative}

Carnap seeks the most comprehensive comparison which agrees with
\emph{every} regular confirmation function, not the numerical ordering
induced by one chosen function. Agreement alone would admit arbitrarily
narrow relations; his second requirement demands that no unanimously
ordered pair be left out. The definition itself nevertheless uses only
logical relations, not numerical measures
\cite[secs.~79--81, pp.~428--440]{Carnap1950Probability}.

\begin{hbdefinition}{def:p15-carnap-comparative}{The purely comparative relation}
Write \(\operatorname{MC}(h,e;h',e')\) when both evidence sentences
are consistent and at least one of the following holds:
\[
 e\models h;\qquad e'\models\neg h';\qquad
 e'\land h'\models e\land h\ \text{ and }\ e\models h\lor e'.
\]
These alternatives are nonexclusive. Define \(\operatorname{Eq}\)
by MC in both directions, \(\operatorname{Gr}\) by MC in the first
direction but not the reverse, and \emph{comparability} by MC in at
least one direction. Maximum confirmation means consistent evidence
entailing the hypothesis; minimum means consistent evidence entailing
its negation. These are Carnap's D81-1, D82-1--3, and D84-1--2 in
modern logical notation.
\end{hbdefinition}

\begin{hbtheorem}{thm:p15-carnap-comparative}{proved=annex:p15.carnap-comparative}{The finite adequacy theorem and its limit boundary}
In any finite language system, with consistent \(e,e'\),
\[
 \operatorname{MC}(h,e;h',e')
 \quad\Longleftrightarrow\quad
 c(h,e)\geq c(h',e')\ \text{for every regular }c.
\]
Thus Eq means equality for every such function; Gr means inequality for
every such function and strict inequality for at least one. Incomparability
means that two regular functions give opposite strict orderings.
For the same evidence, MC is equivalent to
\(e\land h'\models h\); Eq is equivalence of the hypotheses
relative to \(e\). Maximum and minimum correspond to values one and
zero respectively under every regular measure.

If a fitting sequence of finite systems satisfies the comparison
eventually and both confirmation values have limits, the limit values
have the same weak ordering. For nongeneral sentences using only a fixed
finite set of names, the finite characterization also gives the comparison
over all regular fitting extensions to the infinite language. The claim
for arbitrary general sentences requires the additional logical-transfer
and limit-existence conditions discussed below.
\end{hbtheorem}

The proof uses all regular measures, including measures witnessing a
failed comparison; a calculation under \(c^*\) alone cannot establish
this theorem. The relation is in general only a partial ordering of
sentence pairs. For example, under tautological evidence, two independent
atoms \(P(a),Q(a)\) are incomparable: masses concentrated near the
\(P\land\neg Q\) profile or the \(\neg P\land Q\) profile give
opposite predictions. In contrast \(P(a)\lor Q(a)\) is at least as
confirmed as \(P(a)\) on every common consistent evidence.

\begin{hbtheorem}{thm:p15-carnap-comparative-calculus}{proved=annex:p15.carnap-comparative-calculus}{Comparison calculus and predictable observations}
For consistent evidences, MC is reflexive and transitive; Eq is an
equivalence relation and Gr is an irreflexive transitive relation.
Complementation reverses comparison:
\[
 \operatorname{MC}(h,e;h',e')\Longrightarrow
 \operatorname{MC}(\neg h',e';\neg h,e).
\]
If MC compares \((h,e)\) with \((h',e')\) and compares
\((j,e\land h)\) with \((j',e'\land h')\), then it compares
\((h\land j,e)\) with \((h'\land j',e')\), provided the two
conjunctive evidences are consistent. The same conclusion holds for
the crossed comparison of these two factors. Conversely, a product
comparison and a reverse comparison of one positive factor imply the
forward comparison of the other factor. If a fixed left-hand pair
dominates both \((h',e'\land i)\) and \((h',e'\land j)\), it
dominates \((h',e'\land(i\lor j))\), when these evidences are
consistent.

If \(e\land h\models i\), \(e\land h\) is consistent, and
\(e\land i\) is consistent, then for every regular finite \(c\),
\[
 c(h,e\land i)=\frac{c(h,e)}{c(i,e)}\geq c(h,e).
\]
The inequality is strict precisely when \(e\not\models i\).
For two such predictable observations \(i,j\), their posterior
confirmations of the same \(h\) have the reverse ordering of their
prior expectedness \(c(i,e),c(j,e)\). For two hypotheses which
both predict the same observation, their posterior ordering is their
prior ordering.
\end{hbtheorem}

This reconstructs the comparison and composition results of §§83--85,
including the probabilistic content of the Koopman comparisons, without
mistaking a statement valid for each selected \(c\) for a statement
whose premise holds for all \(c\). Carnap explicitly explains this
quantifier distinction in §85, pp.~456--457.

\section{Relevance, quotients, and the signed relevance measure}
\label{sec:p15-carnap-relevance}

\begin{hbdefinition}{def:p15-carnap-relevance}{Carnap's relevance convention}
Additional evidence \(i\) is positive to \(h\) on \(e\) if
\(c(h,e\land i)>c(h,e)\), negative if the inequality is reversed,
and relevant in either case. It is irrelevant if the two defined values
are equal \emph{or} \(e\land i\) is logically contradictory.
The latter is Carnap's explicit D65-1d convention; it assigns no
posterior probability to contradictory evidence. Initial relevance is
relevance on tautological evidence.

When \(c(h,e),c(i,e)>0\), the simple relevance quotient and its
multiple version are
\[
 Q_e(h,i)\defeq\frac{c(h,e\land i)}{c(h,e)},\qquad
 Q_e(i_1,\ldots,i_n)\defeq
 \frac{c(\bigwedge_{r=1}^n i_r,e)}{\prod_{r=1}^n c(i_r,e)}.
\]
All displayed quotients require positive denominators and defined
conditioning values. Independently, wherever its measure values exist,
Carnap's signed relevance measure is
\[
 r(i,h,e)\defeq
 m(e\land i\land h)m(e)-m(e\land h)m(e\land i).
\]
It is not a nonnegative probability measure.
\end{hbdefinition}

\begin{hbtheorem}{thm:p15-carnap-relevance-core}{proved=annex:p15.carnap-relevance-core}{Four-cell determinant and additivity}
Put
\[
 (a,b,c,d)\defeq
 \bigl(m(eih),m(ei\neg h),m(e\neg i h),m(e\neg i\neg h)\bigr),
\]
where juxtaposition abbreviates conjunction. Then \(r=ad-bc\).
When \(m(e),m(ei)>0\),
\[
 c(h,ei)-c(h,e)=\frac{r(i,h,e)}{m(e)m(ei)}.
\]
In a finite regular system, positive, negative, and irrelevant correspond
exactly to \(r>0,r<0,r=0\), including Carnap's contradiction
convention. Also
\[
 r(i,h,e)=r(h,i,e),\quad
 r(\neg i,h,e)=-r(i,h,e),\quad
 r(i,\neg h,e)=-r(i,h,e).
\]
If \(m(e)>0\), \(r=0\) is equivalent to
\(c(ih,e)=c(i,e)c(h,e)\).
For arbitrary \(i,j\),
\[
 r(i\lor j,h,e)=r(i,h,e)+r(j,h,e)-r(ij,h,e).
\]
If \(m(eij)=0\), the last term vanishes. If
\(m(e\neg i\neg j)=0\), the conjunction's relevance instead is
the sum of the two relevances. These identities hold also in the
hypothesis argument by symmetry. In nonfinite systems the algebraic
identities require existence of the displayed measure values; recovering
conditional comparisons requires the positive-denominator conditions.
\end{hbtheorem}

\begin{hbtheorem}{thm:p15-carnap-relevance-quotients}{proved=annex:p15.carnap-relevance-quotients}{Quotients and accumulated evidence}
On their stated domains, \(Q_e(h,i)=Q_e(i,h)\), and
\(Q_e(h,i)=1+r(i,h,e)/(m(eh)m(ei))\). The multiple quotient is
unchanged by permutation of its arguments, and
\[
 c(\bigwedge_r i_r,e)=Q_e(i_1,\ldots,i_n)\prod_r c(i_r,e).
\]
For hypotheses \(h,h'\) and observations \(i_1,\ldots,i_n\),
let \(O_0=c(h,e)/c(h',e)\) and
\(O_r=c(h,ei_r)/c(h',ei_r)\). If every probability used in a
denominator is positive, including the two conditional multiple
quotients, then
\[
 \frac{c(h,e\bigwedge_r i_r)}{c(h',e\bigwedge_r i_r)}
 =\frac{Q_{eh}(i_1,\ldots,i_n)}{Q_{eh'}(i_1,\ldots,i_n)}
   O_0^{1-n}\prod_r O_r.
\]
Thus separate posterior ratios do not determine the accumulated ratio
without the dependence quotients under both hypotheses.
\end{hbtheorem}

These are the Johnson--Keynes quotient methods reported in §66,
pp.~356--360, followed by Carnap's distinct signed-measure method of
§67. The change in confirmation itself is generally asymmetric under
exchanging hypothesis and observation; the signed measure is symmetric
because it includes the appropriate factors of the evidence masses.

\section{Two observations, two hypotheses, and all sign cases}
\label{sec:p15-carnap-relevance-cases}

\begin{hbtheorem}{thm:p15-carnap-relevance-sign-cases}{proved=annex:p15.carnap-relevance-sign-cases}{The four-cell method and the seventy-five situations}
For observations \(i,j\), put
\[
 (x_1,x_2,x_3,x_4)=
 \bigl(r(ij,h,e),r(i\neg j,h,e),r(\neg i j,h,e),
       r(\neg i\neg j,h,e)\bigr).
\]
Then \(\sum x_l=0\), and the relevances of \(i,j,i\lor j\)
are \(x_1+x_2,x_1+x_3,-x_4\). The complete set of possible seven
sign-tuples is obtained by taking these signs for every real zero-sum
four-vector. There are exactly seventy-five distinct tuples: thirty-eight
with no zero among the four cell signs, twenty-four with exactly one
zero, twelve with exactly two zeros, and the all-zero tuple. Every such
tuple is realized by a regular finite state model.

This gives the following complete coarse classification. If the two
observations are both positive, or one is positive and the other
irrelevant, at least one of their conjunction and disjunction is positive;
the other may have any of the three signs. The negative version is
obtained by reversing signs. If one is positive and the other negative,
the conjunction and disjunction may have any of the nine sign
combinations. If both are irrelevant, the conjunction and disjunction
are either both irrelevant or have opposite signs and opposite
relevance measures. Logical exclusion sets \(x_1=0\); logical
exhaustiveness sets \(x_4=0\); implication sets the appropriate
off-diagonal cell to zero. The same complete analysis holds for a single
observation and two hypotheses by interchanging arguments of \(r\).
\end{hbtheorem}

This is a generating form of the complete classification in T69-1--3
and T71-2--3, rather than an appeal to an unspecified sample of cases.
It also handles all deductive special cases in T68-4 and T70-4 by
setting logically impossible cell masses to zero. In an infinite system,
the signs give conditional relevance only under the positive-mass
conditions of T69-3 and T71-3; zero relevance measure alone does not
resolve an almost-logically-false conditioning sentence.

\begin{hbexample}{def:p15-carnap-chess-relevance}{The source's chess-tournament countercases}
On prior evidence \(e\), ten players have equal chances of winning.
Their composition is:
\[
\begin{array}{c|cc}
 &\text{local}&\text{visiting}\\\hline
 \text{junior}&1\text{ man},2\text{ women}&2\text{ men}\\
 \text{senior}&2\text{ men}&3\text{ women}
\end{array}
\]
For the hypothesis that a man wins, prior confirmation is \(1/2\).
Learning that the winner is local gives \(3/5\); learning that the
winner is junior also gives \(3/5\); learning both gives \(1/3\).
For a woman winning, learning visiting or learning senior separately
gives \(3/5>1/2\), but their disjunction gives \(3/7<1/2\).

Equally, the report that a man wins raises both the local and junior
hypotheses from \(1/2\) to \(3/5\), but lowers their conjunction
from \(3/10\) to \(1/5\). The report that a woman wins raises both
the visiting and senior hypotheses from \(1/2\) to \(3/5\), but
lowers their disjunction from \(7/10\) to \(3/5\).
The complementary reports give the sign-reversed countercases.
\end{hbexample}

\section{Disjunctive and conjunctive analysis}
\label{sec:p15-carnap-relevance-analysis}

\begin{hbdefinition}{def:p15-carnap-extreme-complete}{Extreme and complete relevance}
An observation \(i\) is \emph{extremely positive} to \(h\) on
\(e\) if it is positive and no stronger sentence \(j\), with
\(e\land j\models i\), is negative. It is \emph{completely
positive} if it is positive and no weaker sentence \(j\), with
\(e\land i\models j\), is negative. Replace positive/negative
by negative/positive for the corresponding negative definitions.
Extreme irrelevance requires irrelevance of \(i\) and of every such
stronger \(j\); complete irrelevance requires irrelevance of \(i\)
and of every such weaker \(j\). These definitions are meaningful
generally wherever the underlying relevance concepts are defined.
\end{hbdefinition}

\begin{hbtheorem}{thm:p15-carnap-relevance-analysis}{proved=annex:p15.carnap-relevance-analysis}{State and content analyses; full extreme/complete criteria}
In a finite system, let \(z\) range over complete state descriptions.
Then
\[
 r(i,h,e)=\sum_{z\models i}r(z,h,e)
          =\sum_{z\models\neg i}r(\neg z,h,e).
\]
Terms for \(z\not\models e\) vanish. For the others,
\[
 r(z,h,e)=
 \begin{cases}
 m(z)m(e\neg h),&z\models h,\\
 -m(z)m(eh),&z\models\neg h.
 \end{cases}
\]
The negations of states falsifying \(i\) are its \emph{content
elements}; their conjunction is equivalent to \(i\). Stronger
sentences have more such excluded-state content, not more admissible
states.

With the four masses \((a,b,c,d)\) defined above, the necessary and
sufficient conditions are
\[
\begin{array}{c|c}
 \text{extremely positive}&b=0, a>0, d>0\\
 \text{completely positive}&c=0, a>0, d>0\\
 \text{extremely negative}&a=0, b>0, c>0\\
 \text{completely negative}&d=0, b>0, c>0.
\end{array}
\]
Extreme irrelevance means \(ei\), \(eh\), or \(e\neg h\) is
contradictory. Complete irrelevance means \(e\neg i\), \(eh\),
or \(e\neg h\) is contradictory. These stronger classifications
are therefore independent of the particular regular measure chosen.
Extreme positive/negative relevance passes to stronger compatible
observations; complete positive/negative relevance passes to weaker
observations not already entailed by the prior evidence.

For ordinary positive relevance, extreme and complete mean respectively
\(b=0\) and \(c=0\); neither means all four cells are nonempty
and \(ad>bc\); both means \(b=c=0\), \(a,d>0\).
The negative version uses \(a,d\) in place of \(b,c\).
Universal positive relevance over all regular measures holds exactly
when \(a,d\) are positive and at least one of \(b,c\) is zero;
universal negative relevance has the analogous reversed condition.
Universal irrelevance means both products are structurally zero.
Negating the observation interchanges extreme and complete and reverses
the sign; negating the hypothesis reverses the sign without that
interchange. Swapping observation and hypothesis interchanges extreme
and complete for positive relevance, while preserving their negative
forms.
\end{hbtheorem}

These formulas reconstruct both analyses and the full logical cell
classification of §§72--76, including the distinction between an
impossible component and cancellation among nonzero components.
Keynes's stricter irrelevance requires that the observation supply no
new relevant content. In a finite system it is complete irrelevance
with compatible \(ei\); if the prior evidence decides neither
\(h\) nor \(\neg h\), this leaves only observations already
entailed by \(e\). It cannot replace ordinary irrelevance.

For example, in the binary \(c^*\) model after \(n\) positives
and \(n\) negatives, the prediction for another positive is \(1/2\).
An additional positive alone raises it to \((n+2)/(2n+3)\), and
an additional negative alone lowers it to \((n+1)/(2n+3)\); both
together restore \(1/2\). The combined evidence is ordinarily
irrelevant but has relevant content parts. Requiring every weaker
sentence to be strictly positive would also wrongly exclude tautologies
and context-entailed sentences, which are irrelevant.

\begin{hbtheorem}{thm:p15-carnap-infinite-relevance-boundary}{proved=annex:p15.carnap-infinite-relevance-boundary}{An actual infinite-system exception from finite limits}
In the binary \(c^*\) construction, let \(e\) report \(s\)
positive named individuals, let \(h=\forall xP(x)\), and let
\(c\) name a further individual. For \(N>s\),
\[
 c_N^*(h,e)=\frac{s+1}{N+1},\qquad
 c_N^*(h,eP(c))=\frac{s+2}{N+1}.
\]
Both limits are zero, so the finite strict increase disappears in the
limit. Yet \(h\) entails \(P(c)\): its limit conditional
confirmation of \(P(c)\) is one, greater than
\((s+1)/(s+2)\). Thus symmetry of ordinary positive relevance
fails without the source's non-null limiting conditions.

Let \(j=\forall x(x\ne c\to P(x))\). Then
\(m_N^*(j)=1/N\) and \(c_N^*(h,ej)=N/(N+1)\to1\), whereas
\(c_N^*(h,e)\to0\) and \(r_N^*(j,h,e)\to0\).
The positive limiting change is not recovered from the limiting signed
measure. The limiting conditional value is a limit of finite ratios,
not a division by the zero limiting evidence mass.
\end{hbtheorem}

Carnap distinguishes these almost-logically-false cases from actual
contradictions in §§65 and 67. He also discusses an alternative relevance
definition that counts a transition from probability one to logical
entailment as positive even when the number does not rise. We retain
his adopted numerical definition, not that alternative, and do not
silently replace contradictory evidence by evidence of zero limiting mass.

\section{Classificatory confirmation and the Hempel comparison}
\label{sec:p15-carnap-classificatory}

Carnap's purely logical candidate \(C'(h,i,e)\) requires consistent
\(eih\) and \(e\neg i\neg h\), and either
\(eh\models i\) or \(ei\models h\). The extreme/complete
criterion shows that this is precisely strict positive relevance under
every finite regular function. Its quantifier formulation differs from
Gr between posterior
and prior pairs, whose definition uses a universal weak comparison plus
a strict comparison for at least one function. In §86 Carnap regards
\(C'\) as too narrow, not as his finished classificatory explication:
new ordinary instances and disjoint-sample predictions need not satisfy
either entailment. His quantitatively defined candidate is simply
\(C^*(h,i,e)\) when \(c^*(h,ei)>c^*(h,e)\).

The actual criticism in §§87--88 concerns the meaning of
\emph{confirming}, not the soundness of deductive implication.
Equivalence of hypotheses and compatibility with the evidence remain
appropriate; an unconditional entailment-confirmation principle must
exclude a hypothesis already certain on the prior evidence. Consequence,
converse-consequence and conjunction closure do not hold for positive
relevance, as the tournament cases demonstrate. A high posterior threshold
is a different explicandum. For instance an unchanged value of \(3/4\)
exceeds \(1/2\) but supplies no increase at all.

Carnap reports Hempel's development of a formula over the individuals
essentially named in an observation report, and its consequence-closed
confirmation relation. His objection is that it covers all-positive
sample cases but omits many mixed-frequency and additional-neutral-data
cases of ordinary scientific confirmation. He concludes that the purely
nonquantitative classificatory task is still open. This is Carnap's
1950 assessment of Hempel, distinct from the separately allocated
Hempel-original presentation and from the later historical chapter.

\section{A general estimate, not only a frequency rule}
\label{sec:p15-carnap-general-estimation}

Sections~98--100 discuss a general quantitative language, conditional
on a suitable confirmation function already being available. Carnap
expressly says that its construction for the full quantitative language
of physics remains a future task. That assumption cannot serve as a
proof that such a language-wide function exists.

The contrasts in §98 are historical claims about rival methods, not
a description of all present-day statistics. Carnap distinguishes
maximum likelihood, which selects a parameter making the observed
sample most likely, from an evidence-relative confirmation-weighted
estimate. He distinguishes Neyman's repeated-sampling coverage from
confirmation that a particular unknown parameter is in the reported
interval. And he distinguishes Wald's minimization of maximum risk
from minimizing confirmation-weighted loss or utility loss. He credits
these methods' practical achievements while arguing that one chosen
confirmation function would unify the inductive tasks. The substantive
comparison of their outputs, including small-sample criticism of maximum
likelihood, is explicitly deferred to a later volume; the definitions
alone are not that vindication.

\begin{hbdefinition}{def:p15-carnap-c-mean}{The \(c\)-mean estimate and its distribution}
For an unknown quantity \(X=f(u)\), suppose the hypotheses
\(h_j\) saying \(X=v_j\) are mutually exclusive and exhaustive
relative to consistent evidence \(e\). Carnap's estimate is
\[
 \operatorname e_c(f,u,e)\defeq\sum_j v_jc(h_j,e).
\]
The function and its argument identify the unknown quantity, not just
its actual numerical value. We abbreviate the estimate by
\(\widehat X_e\).

For a general real-valued quantitative language, suppose its chosen
confirmation assignments give a probability law \(\mu_e\) for
\(X\). Require \(\int |x|\,d\mu_e(x)<\infty\) and define
\[
 \widehat X_e\defeq\int x\,d\mu_e(x).
\]
Its cumulative confirmation distribution is
\(F_e(t)=\mu_e(({-\infty},t])\); thus an interval \((a,b]\)
has confirmation \(F_e(b)-F_e(a)\), and the integral may be written
\(\int x\,dF_e(x)\). If \(\mu_e\) has a density \(p_e\),
it is \(\int xp_e(x)\,dx\). A density is an additional condition,
not a consequence of every exact-value hypothesis having probability
zero. This is the careful measure-theoretic form of the density and
Stieltjes alternatives on pp.~526--527.
\end{hbdefinition}

\begin{hbtheorem}{thm:p15-carnap-general-mean}{proved=annex:p15.carnap-general-mean}{Well-defined general means and linear transformations}
Suppose a quantitative language is represented by measurable quantities
on a probability space and evidence by an event \(E\) of positive
probability. Its conditional measure
\(P_E(A)=P(A\cap E)/P(E)\) and the distribution of a real measurable
\(X\) under this measure exist. The finite-value definition agrees
with its distributional integral. If \(X_1,\ldots,X_l\) are
integrable under \(P_E\), then
\[
 \widehat{b_0+\sum_j b_jX_j}_e=b_0+\sum_j b_j\widehat X_{j,e}
 \quad(b_j\in\mathbb R).
\]
No independence is required. An estimate need not be a possible value
or a most probable value. A general real-valued law need not have a
finite mean: the density \(1/(\pi(1+x^2))\) has both positive and
negative part integrals infinite, so its mean is undefined.
\end{hbtheorem}

This theorem supplies the general integral mathematics \emph{conditional
on} a chosen law, not a construction of Carnap's missing full scientific
inductive language. Its integral prerequisites and proofs are in
Chapter~\ref{sec:p15-integration-expectation}.

The original lottery example has prizes \(10,2,0\) and confirmations
\(0.01,0.15,0.84\). Its estimate is \(0.4\), although the most
probable prize is zero and a prize of \(0.4\) cannot occur. The
equiprobable values \(3,4,5\) have mean four; changing only the last
value to \(500\) gives mean \(169\), while the median remains
four. Carnap uses these changes to argue that a monetary valuation
should notice the sizes of outcomes, not just their ranking. That is a
decision-theoretic argument with a purpose and utility assumptions, not
a proof that means solve every estimation or decision problem.

\section{Reliability is itself an estimate}
\label{sec:p15-carnap-reliability}

In §102 Carnap considers three explications: confirmation in a
neighborhood of the estimate (or its limiting density); the width of
a neighborhood with a prescribed confirmation such as one half;
and an estimate of the error. He adopts a squared-error version of the
third, because it accounts for remote possible values as well as nearby
ones. A density-based proposal needs a density and an existing limit;
an exact-coverage interval needs to exist. These cannot be presumed
for every discrete or mixed distribution.

\begin{hbdefinition}{def:p15-carnap-estimated-error}{Error, estimated square error, and standard error}
Use Carnap's sign convention \(V\defeq\widehat X_e-X\): error is
positive when the estimate is too high. Its estimated absolute error
is \(\int|V|\,dP_E\). If \(X\) is square integrable, define
\[
 J_c^2(X,e)\defeq\int V^2\,dP_E,\qquad
 J_c(X,e)\defeq\sqrt{J_c^2(X,e)}.
\]
These are the estimated square error and estimated standard error of
D103-1. They are not the actual error, which depends on the unknown
value, nor an externally established quality score for the chosen
confirmation function.
\end{hbdefinition}

\begin{hbtheorem}{thm:p15-carnap-estimation-reliability}{proved=annex:p15.carnap-estimation-reliability}{Error identities and the nonlinear estimation problem}
For integrable \(X\), the estimate of its signed error is zero.
For square-integrable \(X\),
\[
 J_c^2(X,e)=\widehat{X^2}_e-(\widehat X_e)^2,
 \qquad
 \widehat{(a-X)^2}_e=J_c^2(X,e)+(a-\widehat X_e)^2.
\]
Thus the \(c\)-mean uniquely minimizes estimated squared loss over
real constant estimates \(a\). The estimated square error is zero
exactly when \(X\) is equal to its mean almost surely under the
chosen conditional law. In a regular finite state model this means
evidence admits only that value, not merely that it is very probable.

For equiprobable values \(1,2,3\), the mean is two but the mean
square is \(14/3\), not four. The three confirmation vectors
\((1/3,1/3,1/3)\), \((1/5,3/5,1/5)\), and
\((2/5,1/5,2/5)\) all give mean two but estimated square errors
\(2/3,2/5,4/5\), respectively. The lottery above has estimated
square error \(1.44\) and standard error \(1.2\).
\end{hbtheorem}

The difference between \(\widehat{X^2}_e\) and
\((\widehat X_e)^2\) is Carnap's ``paradox of estimation'' in
§100C, pp.~530--533. It is not an algebraic contradiction: these
answer different questions. Treating a weighted mean as an actual
predicted value creates the misleading two decisions. Carnap resolves
the practical problem by estimating the gain or utility of the
\emph{action}, not replacing every unknown intermediate quantity by
its mean before a nonlinear transformation.

\begin{hbwarning}{def:p15-carnap-reliability-boundaries}{Neighborhoods, tails, and self-validation}
A law with probabilities \(0.99,0.005,0.005\) at \(0,-100,100\)
has mean zero and probability \(0.99\) of lying within one unit
of it, but estimated square error \(100\). A law with probabilities
\(0.99,0.01\) at \(0,100\) has mean one; no symmetric closed
interval centered at one has probability exactly \(1/2\).
These are concrete reasons not to assume all of §102's proposed
explications are universally available or interchangeable.

More importantly, §103, pp.~539--540, rejects comparing two
confirmation functions by their own estimated standard errors.
A poorly justified function may assign both an implausible estimate and
an implausibly small error. The error calculation is internal to the
chosen function. It can compare estimates made with that function on
different evidence once the function has other grounds of adequacy;
it cannot supply those grounds by itself.
\end{hbwarning}

\section{Truth frequencies and direct/predictive reliability}
\label{sec:p15-carnap-frequency-reliability}

\begin{hbtheorem}{thm:p15-carnap-frequency-reliability}{proved=annex:p15.carnap-frequency-reliability}{Frequency estimates, finite-population errors, and predictive limits}
For any finite list of sentences \(h_1,\ldots,h_t\), let \(T\)
count its true entries, counting distinct expressions separately even
if logically equivalent. Then
\[
 \widehat T_e=\sum_{j=1}^t c(h_j,e),\qquad
 \widehat{T/t}_e=\frac1t\sum_jc(h_j,e)\quad(t>0).
\]
Logical or inductive dependence does not invalidate these identities.
For a property evaluated on \(t\) new individuals, individual
symmetry makes all the terms equal to the singular prediction \(p\).
The estimated relative frequency is \(p\), independent of sample
size and of which new individuals are selected.
For any such truth count, including overlapping renamed hypotheses,
its estimated square error is
\[
 J_c^2(T/t,e)=\frac1{t^2}\left[
 \sum_j p_j(1-p_j)+2\sum_{j<l}
       \bigl(c(h_jh_l,e)-p_jp_l\bigr)\right],
 \qquad p_j=c(h_j,e).
\]

For direct estimation, suppose evidence specifies \(K\) positives
in a population of \(N\geq2\), and a fixed sample has
\(1\leq t\leq N\) members. Put \(p=K/N\). For its positive
count \(S\), under any regular symmetric measure,
\[
 \widehat S_e=tp,\quad
 J_c^2(S,e)=tp(1-p)\frac{N-t}{N-1},\quad
 J_c^2(S/t,e)=\frac{p(1-p)}t\frac{N-t}{N-1}.
\]
Also \(\widehat{S^2}_e=(tp)^2+J_c^2(S,e)\). For \(N=1\)
the permitted full sample is determined and both errors are zero.
The count mode is governed by
\(q=(t+1)(K+1)/(N+2)\): among feasible counts it is
\(\lfloor q\rfloor\), or the two counts \(q-1,q\) if
\(q\) is an integer and both are feasible. For large \(N\)
with \(t/N\to0\), the variance factor approaches one; it does
not do so for samples covering a substantial fraction of the population.

For predictive estimation under a symmetric conditional law, let all
new indicators have mean \(p\) and distinct-pair joint probability
\(q_2\). Put \(\gamma=q_2-p^2\). Then
\[
 J_c^2(S/t,e)=\frac{p(1-p)}t+\frac{t-1}t\gamma.
\]
If the conditional law fits arbitrarily long exchangeable new samples,
\(\gamma\geq0\) and the estimated square error tends to
\(\gamma\). The limit of the \emph{estimates} is not a proof
that a particular world's empirical relative frequency converges.
For binary \(c^*\) after \(s\) observations with \(s_+\)
positives, \(p=(s_++1)/(s+2)\) and
\(\gamma=p(1-p)/(s+3)\).
\end{hbtheorem}

These are the truth-frequency identity of T104-2, the direct estimates
and reliability quantities of T105-1, and the predictive interpretation
and tentative reliability proposal of §106, including its continuation
on p.~555. The covariance formula is an explicit reconstruction of
that proposal's mathematical content, not a later lambda-family
substitution. For the source's small-population example
\(N=14,K=10,t=7\), the mean count is five and its estimated
square error is \(10/13\); dropping the population correction gives
\(10/7\), nearly twice as large. At \(t=N\) the corrected error
is zero, while that approximation generally is not.

For binary \(c^*\), both no observations and a balanced sample
of one hundred give \(p=1/2\), but their limiting predictive
frequency errors are \(1/12\) and \(1/412\). This illustrates
the distinction between a probability value and the weight of its
evidential support, without turning internal reliability into an external
vindication. The inverse mean for a population whose observed sample
contains \(s_+\) positives is
\([s_++(N-s)p]/N\); inverse and predictive estimation remain
different tasks even when a large remainder makes their answers close.

Section~106B distinguishes a constant sequence of predictive estimates
from a sequence of actual frequencies. Reordering an infinite class
containing infinitely many positives and negatives may change its
frequency limit or destroy it. Individual symmetry concerns their logical
status in the stipulated evidence; it is not an empirical proof about
that ordering or the existence of an actual limit. Section~107 likewise
argues for allowing null evidence and rejects the straight rule's exact
zero/one predictions for new factual cases under the stated regularity
conditions; it expressly defers a full comparative vindication to a later
volume. A hypothetical chosen inductive system and its philosophical
adequacy must not be conflated.

For a general hypothesis, §106C tentatively proposes a finite class
of its renamed versions, with prior-evidence names kept fixed. Symmetry
makes their confirmations equal, so the mean truth-frequency estimate
is the original hypothesis's confirmation. The displayed covariance
sum computes the estimated error of this truth fraction even when
versions overlap or are logically equivalent. Carnap suggests using
this as a reliability measure for the confirmation value. It is not
an error in that value as an analytic calculation, and the proposal
is not reported as a completed external validation criterion.

\begin{proofboundary}{def:p15-carnap-continuation-boundary}{Source and infinite-language standing}
The annex proves the displayed general finite-system comparison and
relevance results, the complete seventy-five-sign construction, the
general conditional-integral estimation results with their integrability
conditions, and the specified infinite-limit countercases. These are not
restricted to a selected lambda family or only to two monadic categories;
the binary \(c^*\) cases are explicitly identified examples.

For arbitrary general sentences in Carnap's infinite language, the
logical transfer used in T81-2a and, in the monadic systems, existence
and effective evaluation of the \(m^*\)-limits are proved in
Chapter~\ref{sec:p15-w02b-p15-029}; existence of the limits of an
arbitrary regular sequence is not asserted. The stated limit comparison assumes the eventual
finite comparison and existing limits. A probability-space representation
assumed for the general estimation theorem does not establish the
original program's full quantitative-language confirmation function.
The Hempel-original and historical allocations remain separate; the
historical scholar copy's 2009 footer versus the publisher's 2011
version is not resolved by this original-source continuation. No
whole-1950-book or whole-historical-chapter closure is claimed.
\end{proofboundary}
